\section{Related Work}
\label{sec:related}

\subsection{Agentic Systems}
\label{subsec:related-agentic}

Capable conversational models~\citep{achiam2023gpt} reframed the LLM as a system component rather than a demo, which opened the question of how to let it act. Tool-augmented language models answered first: Toolformer~\citep{schick2023toolformer} showed an LLM can learn when to call an API, and ToolLLM~\citep{qin2024toolllm} scaled the idea to thousands of real-world tools, turning the model from a text generator into a callable executor. ReAct~\citep{yao2022react} closed that executor into a loop by interleaving chain-of-thought reasoning with environment actions, and the resulting think/act cycle has since become the inner core of every agent that followed. Once one agent worked, splitting roles across many agents within a single inference frame became natural---MetaGPT~\citep{hong2024metagpt}, AutoGen~\citep{wu2024autogen}, ChatDev~\citep{qian2024chatdev}, and CAMEL~\citep{li2023camel} each demonstrated different forms of role specialization---and the design space has by now been mapped at survey scale~\citep{wang2024survey,guo2024large}. This lineage now extends to application-level production agents: Anthropic's Claude assistants~\citep{claude3-anthropic2024} and the open-source OpenClaw~\citep{openclaw-github} run as complete production-grade agentic systems---multi-channel, multi-user services with persistent state and a full harness layer governing routing, gating, and lifecycle. 
\subsection{Self-Evolving Agents}
\label{subsec:related-evolving}

Research on self-evolving agents has developed along two lines. The academic line establishes source-level self-evolution as a working primitive on minimal scaffolds. SICA~\citep{robeyns2025self} shows an agent can edit its own implementation and lift its own SWE-Bench score, settling the feasibility question. The Darwin G\"{o}del Machine~\citep{zhang2025darwin} reframes the same primitive as open-ended search over an archive of agent variants; HyperAgents~\citep{zhang2026hyperagents} pushes that style further by making the meta-procedure itself editable. Meta-Harness~\citep{lee2026meta} isolates a finer signal: exposing past execution traces to the coding-agent proposer strengthens iterative improvement beyond what benchmark-score feedback alone yields. The application-level line extends evolution to deployed agents but bounds its scope to text-mutable artifacts. Hermes Agent~\citep{hermes-github}, paired with the DSPy~\citep{khattab2023dspy} and GEPA~\citep{agrawal2025gepa} optimizers that compile and search over prompts, forms the most ambitious cluster on this side; Capability Evolver~\citep{wang2026procedural}, SkillClaw~\citep{ma2026skillclaw}, GenericAgent~\citep{liang2026genericagent}, and EvoAgentX~\citep{wang2025evoagentx} each occupy a different text-mutable layer---behavioral genes and memory capsules, shared skill corpora, markdown SOPs, prompts and workflow graphs---and none touches the harness. MOSS combines the two: source-level evolution applied to an application-level substrate, with the harness layer included and deployment-failure signals replacing benchmark scores.
